\chapter*{Preface}
\addcontentsline{toc}{chapter}{Preface}
\markboth{Preface}{Preface}

This book is one argument, told once and in full. Across a family of papers,
a single idea has been developed in pieces: that one can compute on values
whose meaning is sealed behind a one-way trapdoor, an untrusted machine
evaluating opaque bit strings through opaque lookup tables, while only a
trusted machine can encode and decode. Each paper carried a piece, the
abstraction, its algebra, its confidentiality, a construction, and each had
to re-establish the common ground before reaching its own result. This
monograph sets the common ground once and follows the idea from its
motivation to its frontier as a single line of reasoning.

\paragraph{Who it is for.}
The intended reader is a researcher or graduate student in security,
cryptography, or information theory who wants the whole picture rather than
one theorem. The mathematics it assumes is light and standard: elementary
probability, the basic vocabulary of information theory (entropy, mutual
information, total-variation distance), and a working familiarity with
cryptographic hash functions. It does not assume prior expertise in oblivious
RAM, fully homomorphic encryption, or simulation-based security; one of the
book's first tasks, in fact, is to say carefully why this framework is none
of those. Readers who know the searchable-encryption literature will find
familiar problems approached from an unfamiliar angle; readers who do not
will find the problems introduced from first principles.

\paragraph{How it is organized.}
The book is in six parts, meant to be read in order, though not at a uniform
pace. Part~I is motivation: it states the paradigm in plain terms and draws
the boundary against the cryptographic tools it is not. Part~II is the
keystone, the cipher map and its four measurable properties, on which
everything later rests. Parts~III and~IV are the theory: the algebra and
type structure of cipher maps, and their confidentiality measured at two
scales that do not reduce to each other. Part~V turns to construction: two ways to
build a cipher map, hash-based and linear-algebraic, the second hiding query
frequency at no per-query cost, and the closure machinery that carries data
structures, keys, and whole programs. Part~VI is the frontier: the open problems the framework
raises and the research program they define. Three appendices collect a
primer on the Bernoulli error model the book leans on, a notation reference,
and a map of where each result's proof lives.

\paragraph{How to read it.}
A reader who wants the argument entire should read straight through; the
density rises deliberately from the motivational Part~I to the technical
core, and relaxes again into the roadmap of Part~VI. A reader who wants the
theory can take Parts~II through~IV and treat Part~V as reference; one who
wants the engineering can read Parts~II and~V and consult the confidentiality
theory as needed. A single small example, an encrypted watchlist, threads
through the book, picking up each property and construction as it arrives, so
that the abstraction always has something concrete attached to it.

\paragraph{Relationship to the papers.}
This is a consolidation, not a replacement. The standalone papers, on the
cipher map abstraction, on its quantitative confidentiality, on its algebraic
types, on the linear-algebraic construction, remain the home of the deeper
proofs, and the book summarizes those results and points to them rather than
reproducing them. What the book adds is the cross-cutting view: the place
where the pieces become one treatment, where a construction's cost is shown
to be a consequence of a property proved three parts earlier, and where the
program's open questions can be seen as a whole.

\paragraph{The thesis in one breath.}
If the book has a single thesis it is this: confidentiality here is
\emph{measured}, not assumed negligible. The framework does not promise that
an adversary learns nothing; it promises that what an adversary learns is a
quantity a designer can compute and trade against space, bandwidth, and
accuracy. That is a narrower guarantee than cryptography's strongest tools
make, and a more honest one for the systems that cannot afford them. The
privacy comes from a one-way hash and a uniform representation, not from
hiding access patterns, and keeping that distinction clear is a discipline
the book holds to throughout.

\paragraph{Acknowledgments.}
The framework grew from a set of notes written in 2023 and 2024, and from the
Bernoulli model of approximate data structures developed alongside it; the
debt to that model, which supplies the error theory the trapdoor side builds
on, runs through every part. This draft is a living manuscript, and the
reader who finds an error, or a place where the argument could be made
cleaner, is doing the book a service by saying so.

\bigskip
\noindent Alexander Towell\\
\today
